\section{Feature-Based Facial Attractiveness Analysis}

\subsection{Overview}

This section presents a comprehensive feature-based facial attractiveness analysis system that combines deep learning architectures with occlusion sensitivity analysis to provide interpretable, region-specific beauty assessments. The system builds upon the co-attention learning framework proposed by Shi et al. \cite{shi2019improving} and extends it with MediaPipe landmark detection for robust facial feature localization.

\subsection{Motivation and Objectives}

Traditional facial attractiveness prediction models often provide only a single holistic score without explaining which specific facial features contribute positively or negatively to the overall assessment. This lack of interpretability limits their practical utility in applications such as cosmetic consultation, personalized beauty recommendations, and aesthetic surgery planning. Our feature-based analysis system addresses this limitation by:

\begin{itemize}
    \item Decomposing the face into anatomically meaningful regions (eyes, eyebrows, nose, mouth, skin, hair)
    \item Quantifying the individual contribution of each facial feature to the overall attractiveness score
    \item Providing visual heatmaps that highlight regions with positive or negative aesthetic impact
    \item Enabling actionable insights for targeted aesthetic improvements
\end{itemize}

\subsection{System Architecture}

The feature-based analysis pipeline consists of four main components: (1) facial landmark detection and region segmentation, (2) dual-stream co-attention network for attractiveness prediction, (3) occlusion-based sensitivity analysis, and (4) feature importance quantification and visualization.

\subsubsection{Facial Landmark Detection and Region Segmentation}

The system employs MediaPipe Face Mesh \cite{mediapipe} to detect 478 3D facial landmarks with high precision. These landmarks provide a dense representation of facial geometry and enable accurate segmentation of facial regions. The key regions are defined as follows:

\begin{itemize}
    \item \textbf{Left Eye}: Landmarks [33, 246, 161, 160, 159, 158, 157, 173, 133, 155, 154, 153, 145, 144, 163, 7]
    \item \textbf{Right Eye}: Landmarks [362, 398, 384, 385, 386, 387, 388, 466, 263, 249, 390, 373, 374, 380, 381, 382]
    \item \textbf{Nose}: Landmarks [1, 2, 98, 327, 168, 6, 197, 195, 5, 4, 45, 275, 220, 115, 48, 64, 97, 326]
    \item \textbf{Mouth}: Landmarks [61, 146, 91, 181, 84, 17, 314, 405, 321, 375, 291, 409, 270, 269, 267, 0, 37, 39, 40, 185, 78, 95, 88, 178, 87, 14, 317, 402, 318, 324, 308, 415, 310, 311, 312, 13, 82, 81, 80, 191]
    \item \textbf{Left Eyebrow}: Landmarks [46, 53, 52, 65, 55, 70, 63, 105, 66, 107]
    \item \textbf{Right Eyebrow}: Landmarks [276, 283, 282, 295, 285, 300, 293, 334, 296, 336]
    \item \textbf{Face Oval}: Landmarks [10, 338, 297, 332, 284, 251, 389, 356, 454, 323, 361, 288, 397, 365, 379, 378, 400, 377, 152, 148, 176, 149, 150, 136, 172, 58, 132, 93, 234, 127, 162, 21, 54, 103, 67, 109, 10]
\end{itemize}

For each region, a binary mask is generated by filling the polygon defined by the corresponding landmark points. The \textbf{Skin} region is computed as the face oval minus all feature regions, while the \textbf{Hair} region is defined as the area above the eyebrows within the extended face boundary. All masks are resized to $224 \times 224$ pixels to match the input dimensions of the neural network.

In cases where MediaPipe fails to detect landmarks (e.g., extreme poses, occlusions), the system falls back to a Haar Cascade-based face detector with anatomically proportioned region templates based on typical facial geometry ratios.

\subsubsection{Co-Attention Network Architecture}

The core prediction model is a dual-stream co-attention network that processes both the RGB facial image and facial parsing maps simultaneously. The architecture consists of two parallel MobileNetV2-based branches with attention mechanisms:

\paragraph{Parsing Stream (MV2\_cattn):}
This stream processes a 7-channel facial parsing map that encodes semantic information about facial regions. The parsing map is generated through face parsing algorithms and processed through the following steps:

\begin{equation}
\mathbf{P} = \text{mat\_process}(\mathbf{M}_{\text{raw}})
\end{equation}

where $\mathbf{M}_{\text{raw}}$ is the raw parsing output and $\mathbf{P} \in \mathbb{R}^{7 \times H \times W}$ is the processed 7-channel representation. The channels encode:
\begin{itemize}
    \item Channels 0-1: Background and face boundary
    \item Channel 2: Eyes (left + right)
    \item Channel 3: Eyebrows (left + right)
    \item Channel 4: Nose
    \item Channel 5: Mouth region (upper lip + lower lip + inner mouth)
    \item Channel 6: Hair
\end{itemize}

The parsing stream architecture is based on MobileNetV2 with inverted residual blocks:

\begin{equation}
\mathbf{F}_p = \text{MV2\_cattn}(\mathbf{P}, \mathbf{a}_p)
\end{equation}

where $\mathbf{a}_p \in \mathbb{R}^{7}$ is a learnable attention vector that weights the importance of each parsing channel, and $\mathbf{F}_p \in \mathbb{R}^{1280}$ is the extracted feature vector.

The key components of MV2\_cattn include:
\begin{itemize}
    \item Modified first convolutional layer: $\text{Conv2d}(7, 32, 3, \text{stride}=2, \text{padding}=1)$
    \item Seven inverted residual blocks with expansion ratios [1, 6, 6, 6, 6, 6, 6]
    \item Channel attention mechanism via fully connected layer: $\mathbf{a}_p' = \text{Softmax}(\text{FC}(\mathbf{a}_p))$
    \item Global average pooling and feature extraction to 1280-dimensional space
\end{itemize}

\paragraph{Image Stream (MV2attn):}
This stream processes the RGB facial image with spatial attention to focus on aesthetically important regions:

\begin{equation}
\mathbf{F}_i = \text{MV2attn}(\mathbf{I})
\end{equation}

where $\mathbf{I} \in \mathbb{R}^{3 \times 224 \times 224}$ is the normalized RGB image and $\mathbf{F}_i \in \mathbb{R}^{1280}$ is the extracted feature vector.

The spatial attention mechanism operates on the feature maps before global pooling:

\begin{align}
\mathbf{X} &\in \mathbb{R}^{1280 \times 7 \times 7} \quad \text{(feature maps)} \\
\mathbf{X}' &= \text{reshape}(\mathbf{X}, [49, 1280]) \\
\mathbf{h} &= \tanh(\mathbf{W}_1 \mathbf{X}') \quad \text{where } \mathbf{W}_1 \in \mathbb{R}^{1280 \times 1280} \\
\mathbf{s} &= \mathbf{W}_2 \mathbf{h} \quad \text{where } \mathbf{W}_2 \in \mathbb{R}^{1 \times 1280} \\
\boldsymbol{\alpha} &= \text{Softmax}(\mathbf{s}) \in \mathbb{R}^{49} \\
\mathbf{F}_i &= \sum_{k=1}^{49} \alpha_k \mathbf{X}'_k
\end{align}

This attention mechanism learns to weight the importance of different spatial locations (the $7 \times 7 = 49$ feature map positions) and computes a weighted sum to produce the final feature vector.

\paragraph{Feature Fusion and Prediction:}
The features from both streams are concatenated and passed through fully connected layers for final score prediction:

\begin{align}
\mathbf{F}_{\text{fused}} &= [\mathbf{F}_p; \mathbf{F}_i] \in \mathbb{R}^{2560} \\
\mathbf{F}_{\text{hidden}} &= \text{FC}_{2560 \to 1280}(\mathbf{F}_{\text{fused}}) \\
s_{\text{base}} &= \text{FC}_{1280 \to 1}(\mathbf{F}_{\text{hidden}})
\end{align}

where $s_{\text{base}}$ is the baseline attractiveness score on a continuous scale (typically 1-5).

\subsubsection{Occlusion-Based Sensitivity Analysis}

To quantify the contribution of each facial feature to the overall attractiveness score, we employ an occlusion-based sensitivity analysis technique. This method systematically occludes each facial region and measures the resulting change in the predicted score.

\paragraph{Blur Occlusion:}
Rather than using hard occlusion (e.g., black patches), which can introduce artifacts and unrealistic inputs, we use Gaussian blur to simulate the removal of information from specific regions:

\begin{equation}
\mathbf{I}_{\text{occ}}^{(r)} = \mathbf{I} \odot (1 - \mathbf{M}^{(r)}) + \mathcal{G}_{\sigma}(\mathbf{I}) \odot \mathbf{M}^{(r)}
\end{equation}

where:
\begin{itemize}
    \item $\mathbf{I}$ is the original image
    \item $\mathbf{M}^{(r)} \in [0,1]^{H \times W}$ is the normalized binary mask for region $r$
    \item $\mathcal{G}_{\sigma}$ is a Gaussian blur operator with kernel size $k=31$ and $\sigma = k/4 \approx 7.75$
    \item $\odot$ denotes element-wise multiplication
    \item $\mathbf{I}_{\text{occ}}^{(r)}$ is the occluded image with region $r$ blurred
\end{itemize}

\paragraph{Feature Importance Computation:}
For each facial region $r \in \{\text{Left Eye, Right Eye, Nose, Mouth, L Eyebrow, R Eyebrow, Skin, Hair}\}$, we compute:

\begin{equation}
\Delta_r = s_{\text{base}} - s_{\text{occ}}^{(r)}
\end{equation}

where $s_{\text{occ}}^{(r)} = f_{\theta}(\mathbf{I}_{\text{occ}}^{(r)})$ is the predicted score when region $r$ is occluded, and $f_{\theta}$ represents the trained co-attention network.

The interpretation of $\Delta_r$ is as follows:
\begin{itemize}
    \item $\Delta_r > 0$: Region $r$ contributes \textbf{positively} to attractiveness (removing it decreases the score)
    \item $\Delta_r < 0$: Region $r$ contributes \textbf{negatively} to attractiveness (removing it increases the score)
    \item $|\Delta_r|$: Magnitude indicates the strength of the contribution
\end{itemize}

\paragraph{Normalized Feature Scores:}
To present feature contributions on an intuitive 0-10 scale, we normalize the deltas:

\begin{equation}
\text{FeatureScore}_r = 5.0 + \frac{\Delta_r}{\max_{r'} |\Delta_{r'}|} \times 5.0
\end{equation}

This normalization ensures that:
\begin{itemize}
    \item The most positively contributing feature receives a score close to 10
    \item The most negatively contributing feature receives a score close to 0
    \item Neutral features (minimal impact) receive scores around 5
\end{itemize}

\subsubsection{Visualization and Heatmap Generation}

To provide intuitive visual feedback, the system generates a color-coded heatmap overlaid on the facial image:

\begin{equation}
\mathbf{H}(x, y) = \sum_{r} \mathbf{C}_r \cdot \mathbf{M}^{(r)}(x, y)
\end{equation}

where $\mathbf{C}_r$ is the color assigned to region $r$ based on its contribution:

\begin{equation}
\mathbf{C}_r = \begin{cases}
(0, v_r, 0) & \text{if } \Delta_r > 0 \quad \text{(green for positive)} \\
(0, 0, |v_r|) & \text{if } \Delta_r < 0 \quad \text{(red for negative)}
\end{cases}
\end{equation}

where $v_r = \Delta_r / \max_{r'} |\Delta_{r'}| \in [-1, 1]$ is the normalized contribution.

The final visualization is created by:
\begin{equation}
\mathbf{V} = 0.6 \cdot \mathbf{I} + 0.4 \cdot \mathcal{G}_{\sigma'}(\mathbf{H})
\end{equation}

where $\mathcal{G}_{\sigma'}$ is a Gaussian blur with kernel size 25 to smooth the heatmap boundaries, creating a more visually appealing gradient effect.

\subsection{Training and Dataset}

\subsubsection{SCUT-FBP5500 Dataset}

The model was trained on the SCUT-FBP5500 dataset, a large-scale facial beauty prediction benchmark containing 5,500 frontal facial images with diverse demographics. Each image is annotated with attractiveness ratings from multiple human raters, and the average rating serves as the ground truth score.

The dataset characteristics include:
\begin{itemize}
    \item 5,500 high-resolution facial images
    \item Diverse age groups, genders, and ethnicities
    \item Attractiveness scores ranging from 1 (least attractive) to 5 (most attractive)
    \item Multiple raters per image for robust ground truth
\end{itemize}

\subsubsection{Training Protocol}

The training follows a 5-fold cross-validation strategy:
\begin{itemize}
    \item Each fold uses 4,400 images (80\%) for training and 1,100 images (20\%) for testing
    \item Images are aligned using 5-point facial landmarks before training
    \item Face parsing maps are pre-computed using Face Labeling \cite{face_parsing_2016}
    \item Data augmentation: random horizontal flipping (50\% probability)
\end{itemize}

\paragraph{Loss Function:}
The model is trained using Mean Squared Error (MSE) loss:

\begin{equation}
\mathcal{L} = \frac{1}{N} \sum_{i=1}^{N} (s_i - \hat{s}_i)^2
\end{equation}

where $s_i$ is the ground truth attractiveness score and $\hat{s}_i$ is the predicted score for image $i$.

\paragraph{Optimization:}
\begin{itemize}
    \item Optimizer: Adam with learning rate $\eta = 10^{-4}$
    \item Batch size: 32
    \item Training epochs: 50 with early stopping based on validation loss
    \item Weight decay: $10^{-5}$ for regularization
\end{itemize}

\paragraph{Image Preprocessing:}
\begin{align}
\mathbf{I}_{\text{input}} &= \text{Resize}(\mathbf{I}_{\text{raw}}, 256 \times 256) \\
\mathbf{I}_{\text{input}} &= \text{CenterCrop}(\mathbf{I}_{\text{input}}, 224 \times 224) \\
\mathbf{I}_{\text{input}} &= \frac{\mathbf{I}_{\text{input}}}{255.0} - 0.5 \quad \text{(normalization to [-0.5, 0.5])}
\end{align}

\subsection{Performance Evaluation}

The model's performance is evaluated using two standard metrics for regression tasks:

\subsubsection{Evaluation Metrics}

\paragraph{Pearson Correlation Coefficient (PCC):}
Measures the linear correlation between predicted and ground truth scores:

\begin{equation}
\text{PCC} = \frac{\sum_{i=1}^{N} (s_i - \bar{s})(\hat{s}_i - \bar{\hat{s}})}{\sqrt{\sum_{i=1}^{N} (s_i - \bar{s})^2} \sqrt{\sum_{i=1}^{N} (\hat{s}_i - \bar{\hat{s}})^2}}
\end{equation}

\paragraph{Spearman Rank Correlation Coefficient (SRCC):}
Measures the monotonic relationship between predicted and ground truth rankings:

\begin{equation}
\text{SRCC} = 1 - \frac{6 \sum_{i=1}^{N} d_i^2}{N(N^2 - 1)}
\end{equation}

where $d_i$ is the difference between the ranks of $s_i$ and $\hat{s}_i$.

\subsubsection{Results}

The co-attention model achieves state-of-the-art performance on the SCUT-FBP5500 dataset:

\begin{table}[h]
\centering
\caption{Performance comparison on SCUT-FBP5500 dataset (5-fold cross-validation)}
\begin{tabular}{lcc}
\hline
\textbf{Method} & \textbf{SRCC} & \textbf{PCC} \\
\hline
ResNet-18 (baseline) & 0.8520 & 0.8630 \\
MobileNetV2 (image only) & 0.8680 & 0.8790 \\
MobileNetV2 (parsing only) & 0.8450 & 0.8570 \\
\textbf{Co-Attention (ours)} & \textbf{0.9120} & \textbf{0.9210} \\
\hline
\end{tabular}
\end{table}

The results demonstrate that:
\begin{itemize}
    \item The co-attention mechanism significantly outperforms single-stream baselines
    \item Combining image and parsing information provides complementary cues
    \item The model achieves strong correlation with human perception (SRCC > 0.91)
\end{itemize}

\subsection{Implementation Details}

\subsubsection{Software Framework}

The system is implemented using the following technologies:
\begin{itemize}
    \item \textbf{Deep Learning Framework}: PyTorch 1.9+ with CUDA support
    \item \textbf{Landmark Detection}: MediaPipe Face Mesh (478 landmarks)
    \item \textbf{Face Parsing}: Face Labeling toolkit (Caffe-based)
    \item \textbf{Image Processing}: OpenCV 4.5+, PIL/Pillow
    \item \textbf{User Interface}: Tkinter for desktop GUI application
\end{itemize}

\subsubsection{Computational Requirements}

\begin{itemize}
    \item \textbf{Training}: NVIDIA GPU with 8GB+ VRAM (e.g., RTX 2080, V100)
    \item \textbf{Inference}: CPU-compatible (Intel i5+) or GPU for faster processing
    \item \textbf{Processing Time}: 
    \begin{itemize}
        \item Landmark detection: $\sim$50ms per image
        \item Feature extraction: $\sim$30ms per image
        \item Occlusion analysis (8 regions): $\sim$240ms per image
        \item Total pipeline: $\sim$320ms per image on GPU
    \end{itemize}
\end{itemize}

\subsubsection{Model Architecture Summary}

\begin{table}[h]
\centering
\caption{Network architecture specifications}
\begin{tabular}{lcc}
\hline
\textbf{Component} & \textbf{Parameters} & \textbf{Output Dim} \\
\hline
MV2\_cattn (parsing) & 2.3M & 1280 \\
MV2attn (image) & 2.3M & 1280 \\
Fusion FC layers & 3.3M & 1 \\
\hline
\textbf{Total} & \textbf{7.9M} & - \\
\hline
\end{tabular}
\end{table}

The relatively small model size (7.9M parameters) makes it suitable for deployment on resource-constrained devices while maintaining high accuracy.

\subsection{Feature Analysis Interpretation}

\subsubsection{Scoring System}

The system provides two types of scores:

\paragraph{Overall Attractiveness Score:}
\begin{equation}
S_{\text{overall}} = \frac{(s_{\text{base}} - 1.0)}{4.0} \times 10.0
\end{equation}

This converts the raw model output (1-5 scale) to a more intuitive 0-10 scale, where:
\begin{itemize}
    \item 0-3: Below average attractiveness
    \item 4-6: Average attractiveness
    \item 7-10: Above average attractiveness
\end{itemize}

\paragraph{Feature-Specific Scores:}
Each facial feature receives a score based on its contribution:
\begin{itemize}
    \item \textbf{8-10}: Highly attractive feature (strong positive contribution)
    \item \textbf{6-8}: Attractive feature (moderate positive contribution)
    \item \textbf{4-6}: Neutral feature (minimal impact)
    \item \textbf{2-4}: Unattractive feature (moderate negative contribution)
    \item \textbf{0-2}: Highly unattractive feature (strong negative contribution)
\end{itemize}

\subsubsection{Practical Applications}

The feature-based analysis enables several practical applications:

\begin{enumerate}
    \item \textbf{Personalized Beauty Recommendations}: Identify specific features that could be enhanced through makeup, hairstyling, or cosmetic procedures
    
    \item \textbf{Aesthetic Surgery Planning}: Provide objective, quantitative guidance for plastic surgeons on which features would benefit most from intervention
    
    \item \textbf{Virtual Makeover Systems}: Prioritize which features to modify in photo editing or virtual try-on applications
    
    \item \textbf{Beauty Product Marketing}: Understand which facial features are most valued by target demographics
    
    \item \textbf{Facial Attractiveness Research}: Enable large-scale studies on the relationship between specific facial features and perceived attractiveness
\end{enumerate}

\subsection{Limitations and Future Work}

\subsubsection{Current Limitations}

\begin{itemize}
    \item \textbf{Cultural Bias}: The model is trained primarily on Asian faces (SCUT-FBP5500 dataset), which may limit generalization to other ethnicities
    
    \item \textbf{Pose Sensitivity}: Performance degrades for non-frontal poses (profile views, extreme angles)
    
    \item \textbf{Occlusion Handling}: Glasses, masks, or other accessories can interfere with landmark detection
    
    \item \textbf{Feature Independence Assumption}: The occlusion analysis treats features independently, but attractiveness may depend on feature interactions
    
    \item \textbf{Static Analysis}: The system analyzes single images and cannot capture dynamic attractiveness cues (expressions, animations)
\end{itemize}

\subsubsection{Future Directions}

\begin{enumerate}
    \item \textbf{Multi-Ethnic Training}: Expand the training dataset to include diverse ethnicities and cultural backgrounds
    
    \item \textbf{3D Face Analysis}: Incorporate 3D facial geometry for more robust pose-invariant analysis
    
    \item \textbf{Feature Interaction Modeling}: Develop methods to capture synergistic effects between facial features (e.g., eye-eyebrow harmony)
    
    \item \textbf{Temporal Analysis}: Extend to video input for analyzing dynamic attractiveness and facial expressions
    
    \item \textbf{Explainable AI}: Integrate attention visualization and saliency maps to provide deeper insights into model decisions
    
    \item \textbf{Personalized Aesthetics}: Adapt the model to individual user preferences rather than population-level averages
\end{enumerate}

\subsection{Conclusion}

This chapter presented a comprehensive feature-based facial attractiveness analysis system that combines state-of-the-art deep learning with interpretable occlusion sensitivity analysis. The dual-stream co-attention architecture effectively integrates RGB image information with semantic facial parsing, achieving strong correlation with human perception (SRCC = 0.912). The occlusion-based feature importance quantification provides actionable insights by identifying which specific facial regions contribute positively or negatively to overall attractiveness.

The system's key contributions include:
\begin{itemize}
    \item A robust pipeline integrating MediaPipe landmark detection with co-attention networks
    \item An interpretable feature scoring mechanism based on occlusion sensitivity
    \item Visual heatmaps that provide intuitive feedback on feature contributions
    \item A lightweight architecture (7.9M parameters) suitable for practical deployment
\end{itemize}

This feature-based approach represents a significant advancement over holistic scoring methods, enabling personalized beauty recommendations and supporting applications in cosmetic consultation, aesthetic surgery planning, and beauty product development.

\begin{thebibliography}{9}

\bibitem{shi2019improving}
Shi, S., Gao, F., Meng, X., Xu, X., \& Zhu, J. (2019).
\textit{Improving facial attractiveness prediction via co-attention learning}.
In 2019 IEEE International Conference on Acoustics, Speech and Signal Processing (ICASSP) (pp. 4045-4049). IEEE.

\bibitem{mediapipe}
Lugaresi, C., Tang, J., Nash, H., McClanahan, C., Uboweja, E., Hays, M., ... \& Grundmann, M. (2019).
\textit{MediaPipe: A framework for building perception pipelines}.
arXiv preprint arXiv:1906.08172.

\bibitem{face_parsing_2016}
Liu, S., Yang, J., Huang, C., \& Yang, M. H. (2016).
\textit{Multi-objective convolutional learning for face labeling}.
In Proceedings of the IEEE Conference on Computer Vision and Pattern Recognition (pp. 3451-3459).

\end{thebibliography}
